\section{A limitation of neuronwise cloning, even at polynomial path weight}
\label{sec:neuronwise-cloning-barrier}

A circuit's absolute path weight is the sum, over its input-to-output
paths, of the products of absolute edge weights. Even a circuit with
polynomial absolute path weight can have a
quasipolynomial minimum product of layer Frobenius norms under \emph{all}
neuronwise cloning choices. The statement below concerns an exact
representation-preserving method; it is not a lower bound on the minimum
energy of the represented Boolean function, which may have a different
and much cheaper realization.

\begin{proposition}[Complete-tree cloning obstruction]
Let $H\ge1$, $d=2^H$, and consider a complete binary computation tree with
$H$ hidden layers. Layer $s$ contains $n_s=2^{H-s}$ computational neurons.
Each neuron has two computational incoming edges of absolute weight $a>0$;
the first layer takes its two inputs from distinct Boolean coordinates.
The final scalar readout has absolute weight $\beta>0$.
Arbitrary additional bias-supply neurons and edges are permitted.

Replace every computational neuron by any positive integer number of
identical copies, divide its outgoing edges by the number of its copies,
and fill unused positions with zero neurons to obtain common hidden width
$N$. Write $E_1,\ldots,E_{H+1}$ for the weighted squared layer norms of the
resulting network. Then
\[
 \prod_{s=1}^{H+1}E_s
 \ge d\beta^2(2a^2)^H 2^{H(H-1)/2}.
\]
Consequently its Gaussian energy satisfies
\[
 \kappa_N\ge (H+1)
 \left[d\beta^2(2a^2)^H 2^{H(H-1)/2}\right]^{1/(H+1)}.
\]
The computational input-to-output absolute path sum is merely
$\beta(2a)^H$. For fixed $a,\beta>0$, this is polynomial in $d$, whereas
the lower bound on the product of layer energies is
$\exp(\Omega((\log d)^2))$.
\end{proposition}

\begin{proof}
For a computational node $v$, let $r_v>0$ denote its number of copies
divided by $N$. Discard all non-computational edges when lower-bounding
layer energies. The remaining contributions are
\[
 F_1=2da^2\sum_{v\text{ at level }1}r_v,
 \qquad
 F_s=a^2\sum_{v\text{ at level }s}
       \sum_{u\text{ a child of }v}\frac{r_v}{r_u}
       \quad(2\le s\le H),
 \qquad
 F_{H+1}=\frac{\beta^2}{r_{\rm root}}.
\]
Each $F_s$ is a convex function of the log multiplicities
$y_v=\log r_v$, being a sum of exponentials of affine functions.
Average $y$ over the finite automorphism group of the complete binary
tree. The group preserves every $F_s$ and acts transitively on each level,
so Jensen's inequality decreases every $F_s$ simultaneously and makes
$r_v=r_s$ constant on level $s$. This averaging is used only for a lower
bound and need not give integer clone counts. The width constraints also
remain valid, by convexity of the exponential, although that fact is not
needed for the lower bound.

At these averaged log multiplicities,
\[
 F_1=2da^2 n_1r_1,\quad
 F_s=2a^2 n_s\frac{r_s}{r_{s-1}},\quad
 F_{H+1}=\frac{\beta^2}{r_H}.
\]
Their product telescopes to
$d\beta^2(2a^2)^H\prod_{s=1}^H n_s$.
Since $E_s\ge F_s$ before averaging and every $F_s$ decreases on
averaging, this lower-bounds the original product of layer energies.
Now $\prod_s n_s=2^{H(H-1)/2}$, and arithmetic--geometric mean gives the
energy bound. There are $2^H$ computational input-to-output paths, each
with absolute weight $\beta a^H$, proving the last assertion.
\end{proof}

In particular, allowing nonuniform clone counts within a layer does not
turn a polynomial absolute path bound into the polynomial energy-product
bound needed by the unrestricted Gaussian-distance theorem. A proof of
that implication would need a change of representation beyond exact
cloning, or a different posterior argument.
